\chapter*{ABSTRACT}
\thispagestyle{plain}

\vspace{0.5cm}

Blind and Low Vision (BLV) individuals face persistent challenges in navigating environments and interacting with physical objects due to the absence of visual information. While traditional mobility aids and existing assistive technologies provide partial support, most solutions focus on isolated tasks such as obstacle detection or text recognition and lack integrated, real-time environmental understanding. This project presents Echora, an AI-powered sensory substitution system designed to enhance independence, safety, and situational awareness for BLV users through multimodal perception and feedback.


Echora combines wearable sensing hardware with smartphone-based artificial intelligence to translate spatial and semantic information into intuitive audio and haptic cues. The system integrates depth sensing, computer vision and inertial measurement to support both macro-navigation (environmental awareness and obstacle avoidance) and micro-interaction (precise hand guidance and object interaction). By employing real-time multimodal sensor fusion and adaptive mode switching, Echora delivers context-aware feedback while minimizing cognitive load on the user.

All high-level processing and AI inference are performed locally on the mobile device, ensuring low latency, privacy preservation, and independence from cloud connectivity. The wearable components are limited to sensing and feedback, enabling a lightweight, modular, and scalable system architecture. Key features include continuous spatial audio navigation, haptic collision alerts, object and text recognition, and hand-guided interaction.

The proposed system is designed to complement, rather than replace, existing mobility aids such as white canes or guide dogs. By unifying navigation, interaction, and perception within a single adaptive framework, Echora demonstrates the potential of AI-driven sensory substitution systems to significantly improve autonomy and quality of life for Blind and Low Vision users.


